% !PRINTED 159
\weberpage{159}

\begin{center}
\S.\,49.

\gesperrt{Elimination. Theorem of Bezout.}
\end{center}

If in the two functions

\weq{1}{
\begin{aligned}
f(x) &= a_0 x^n + a_1 x^{n-1} + \cdots + a_n \\
\varphi(x) &= b_0 x^m + b_1 x^{m-1} + \cdots + b_m
\end{aligned}
}

the coefficients $a$, $b$ are themselves again integral rational functions
of a quantity $y$, then the resultant will also be an integral
rational function of $y$, which we shall denote by $F(y)$.
The formation of this equation is the elimination of
$x$ from the two equations (1).

The equation

\weq{2}{F(y) = 0}

has as roots all those values of $y$ for which the two
equations

\weq{3}{f(x) = 0, \qquad \varphi(x) = 0}

have a common root. If the condition (2) is satisfied,
one finds the values of $x$ that satisfy the two equations
(3) by seeking the greatest common divisor
of $f(x)$ and $\varphi(x)$. This common divisor
is of the first degree, and thus determines $x$ rationally, when
only one such common root is present; in the other
case it will be of correspondingly higher degree, and the
determination of $x$ still requires the solution of an equation
of higher degree.

If to each root of the resultant there belongs only one common
root $x$, then the number of pairs of values $x$, $y$ that
satisfy the equations (3) is equal to the degree of the resultant
with respect to $y$. Thus if the coefficients $a$ are of degree $\nu$, $b$
of degree $\mu$ with respect to $y$, then by the determination of degree
of the preceding paragraph the degree of the resultant is

% [source-original mathematical slip, confirmed 2026-08-25 against scan p. 184 / printed p. 159]
% The printed expression n*nu + m*mu reverses the bidegrees; the consistent degree is
% m*nu + n*mu. The facing translation preserves the printed formula verbatim.
\weq{5}{n\nu + m\mu;} % [sic] equation numbering jumps from (3) to (5) in print; no (4) occurs here

so great, then, is the number of common pairs of roots.
This number can, however, still be modified by the fact that
the coefficients may possibly be so constituted that
%CONT>>
